% Technical brief with a cover page, callouts, a diagram and a code sample.
% Tectonic reads UTF-8 directly. Replace each sample passage and comment.
\documentclass[11pt,a4paper]{article}

\usepackage[{{language}}]{babel}
\usepackage[T1]{fontenc}
\usepackage{newpxtext}
\usepackage{microtype}
\usepackage[a4paper,margin=2.5cm,headheight=14pt,headsep=12pt,footskip=14pt]{geometry}
\usepackage{amsmath,amssymb}
\usepackage{xcolor}
\usepackage{enumitem}
\usepackage{fancyhdr}
\usepackage{titlesec}
\usepackage{tcolorbox}
\usepackage{booktabs}
\usepackage{float}
\usepackage{hyperref}

% A restrained palette keeps callouts distinct without overpowering the page.
\definecolor{briefink}{HTML}{292927}
\definecolor{briefaccent}{HTML}{8A493D}
\definecolor{briefwarning}{HTML}{A66B2B}
\hypersetup{colorlinks=true,linkcolor=briefaccent,urlcolor=briefaccent,
  citecolor=briefaccent,breaklinks=true,pdftitle={{{title}}},pdfauthor={{{author}}}}

\pagestyle{fancy}
\fancyhf{}
\lhead{{{course}}}
\rhead{{{title}}}
\cfoot{\thepage}
\renewcommand{\headrulewidth}{0.5pt}
\renewcommand{\headrule}{\hbox to\headwidth{\color{briefaccent}\leaders\hrule height \headrulewidth\hfill}}

\titleformat{\section}{\normalfont\Large\bfseries\color{briefink}}
{\textcolor{briefaccent}{\thesection}}{0.75em}{}
\titlespacing*{\section}{0pt}{1.35\baselineskip}{0.45\baselineskip}
\titleformat{\subsection}{\normalfont\large\bfseries\color{briefink}}
{\thesubsection}{0.7em}{}
\titlespacing*{\subsection}{0pt}{1.1\baselineskip}{0.35\baselineskip}
\setlength{\parindent}{0pt}
\setlength{\parskip}{0.5\baselineskip}
\setlength{\emergencystretch}{2em}
\setlist{leftmargin=1.4em,itemsep=2pt,topsep=3pt}

\newtcolorbox{noteBox}{colback=gray!4,colframe=gray!40,
  fonttitle=\bfseries,coltitle=briefink,boxrule=0.5pt,arc=1mm,
  left=6pt,right=6pt,top=5pt,bottom=5pt}
\newtcolorbox{keyBox}{colback=briefaccent!5,colframe=briefaccent!60,
  fonttitle=\bfseries,title=Key Point,boxrule=0.6pt,arc=1mm,
  left=6pt,right=6pt,top=5pt,bottom=5pt}
\newtcolorbox{cautionBox}{colback=briefwarning!6,colframe=briefwarning!70,
  fonttitle=\bfseries,title=Please Note,boxrule=0.6pt,arc=1mm,
  left=6pt,right=6pt,top=5pt,bottom=5pt}

\begin{document}

  % Cover page: keep this block to one page.
  \begin{titlepage}
    \centering
    {\scshape\large {{kicker}} \par}
    \vspace{0.9cm}
    {\huge\bfseries\color{briefink} {{title}} \par}
    \vspace{0.4cm}
    {\Large {{subtitle}} \par}
    \vspace{1.3cm}
    \begin{tcolorbox}[colback=briefaccent!4,colframe=briefaccent!50,
        width=0.88\textwidth,boxrule=0.6pt,arc=1mm]
      \begin{tabular}{@{}ll@{}}
        \textbf{Author:} & {{author}} \\
        \textbf{Course or project:} & {{course}} \\
        \textbf{Date:} & {{date}} \\
        \textbf{Repository:} & {\small\url{{{repository}}}} \\
      \end{tabular}
    \end{tcolorbox}
    \vfill
  \end{titlepage}

  % Replace with the document's purpose and key result.
  \section{Summary}
  \begin{keyBox}
    This brief summarises the current status, the proposed solution and the
    main result for a quick first read.
  \end{keyBox}

  % Replace with the request, the important constraint and what is delivered.
  \section{Context}
  The project needs a concise status update: what was requested, which
  constraint matters and what this version delivers.
  \begin{noteBox}
    Scope of this brief: work completed up to the date shown on the cover.
  \end{noteBox}

  % Replace with the solution steps, diagram and a short code sample.
  \section{Solution}
  The workflow has three steps: prepare the data, apply the main rule and
  present the result.

  \begin{mermaid}[style=editorial,width=0.82\linewidth,
      caption={Overview of the solution workflow},pos=H]
    flowchart LR
    A[Prepare data] --> B[Apply rule] --> C[Present result]
  \end{mermaid}

  \begin{code}[lang=python,numbers=true]
def summarize(items):
    total = sum(items)
    count = len(items)
    mean = total / count if count else 0
    return total, mean
  \end{code}

  % Replace the sample values and the caution with your own results.
  \section{Results}
  The table summarises three illustrative measurements.
  \begin{table}[H]
    \centering
    \caption{Illustrative measurements}
    \begin{tabular}{@{}lcc@{}}
      \toprule
      Case & Before & After \\
      \midrule
      A & 12 & 8 \\
      B & 20 & 11 \\
      C & 15 & 9 \\
      \bottomrule
    \end{tabular}
  \end{table}
  \begin{cautionBox}
    These values are illustrative; verify all figures against the final data.
  \end{cautionBox}

  % Replace with a brief conclusion and the open tasks.
  \section{Conclusion and Next Steps}
  The work meets the stated objective and leaves a clear path for the next
  iteration.
  \begin{itemize}[label=$\square$]
    \item Review the draft with the team.
    \item Add the final data.
    \item Prepare the next release.
  \end{itemize}

  % Replace these generic references with the sources you used.
  \begin{thebibliography}{9}
    \bibitem{guide} Sample Author, \textit{A Short Guide to Technical Briefs}, Sample Press, 2024.
    \bibitem{writing} Sample Author, \textit{Writing with Clarity}, Sample Press, 2023.
  \end{thebibliography}

\end{document}
